In this work, we investigated the ``Judgment Effect'' in large language models, demonstrating that explicitly judgmental and skeptical prompt framing acts as a catalyst for deeper reasoning. Through experiments on the \raio and \causalt datasets using \gptmini, we showed that high-pressure evaluative prompts monotonically increase reasoning trace lengths by up to 22.0\%. Furthermore, on complex social judgment tasks, this increased effort yields an absolute accuracy improvement of 6.6\% over neutral baselines, with the model successfully resisting sycophancy.

Our findings emphasize that users and developers can actively harness conversational tone to regulate an LLM's analytical rigor. Rather than viewing user skepticism purely as a vulnerability that induces deference, it can be repurposed as a structural tool to enforce exhaustive edge-case analysis. Future work should explore the impact of the Judgment Effect on harder logical benchmarks (such as MATH or complex causal tasks) and investigate the threshold at which intense scrutiny shifts from beneficial facilitation to detrimental sycophancy across different state-of-the-art architectures.